\begin{toappendix}
\section{Provenance and Probability in Different Systems}
\label{sec:practical}


\begin{table*}
	\centering
	\caption{Provenance derivation
		for GProM on $\qex$}\label{tab:gprom_query2}
	%\vspace{-1.2em}
	\resizebox{\textwidth}{!}{
		\begin{tabular}{lllllllll}
			\toprule
			\textbf{city} & \textbf{prov\_personnel\_id} & \textbf{prov\_personnel\_name} & \textbf{prov\_personnel\_position} & \textbf{prov\_personnel\_city} & \textbf{prov\_personnel\_1\_id} & \textbf{prov\_personnel\_1\_name} & \textbf{prov\_personnel\_1\_position} & \textbf{prov\_personnel\_1\_city} \\
			\midrule
			\revision{Nairobi}      & 1                            &
      \revision{Juma}                           & Director                           & \revision{Nairobi}                       & 2                               & Paul                              & Janitor                               & \revision{Nairobi}                          \\
      Paris         & 3                            & \revision{David}                           & Analyst                            & Paris                          & 6                               & Nancy                             & HR                                    & Paris                             \\
      Paris         & 3                            & \revision{David}
                    & Analyst                            & Paris
                    & 5                               & \revision{Aaheli}                          & Double agent                          & Paris                             \\
			\revision{Beijing}        & 4                            & Ellen
                                & Field agent                        &
      \revision{Beijing}                         & 7
                                                 & \revision{Jing}                             & Analyst                               & \revision{Beijing}                            \\
      Paris         & 5                            & \revision{Aaheli}                       & Double agent                       & Paris                          & 6                               & Nancy                             & HR                                    & Paris                             \\
			\bottomrule
		\end{tabular}}
\end{table*}
\begin{table}
	\centering
	\footnotesize
	\caption{Provenance derivation for ProvSQL on~$\qex$}\label{tab:provsql_query2}
	%\vspace{-1.2em}
	\begin{tabular}{lll}
		\toprule
		\textbf{city} & \textbf{why}                                                  & \textbf{provsql} \\
		\midrule
    \revision{Nairobi}      & \{"\{\revision{Juma},Paul\}"\}                                           & d1b22232-\dots   \\
    Paris         &
    \{"\{\revision{David},\revision{Aaheli}\}","\{\revision{Aaheli},Nancy\}","\{\revision{David},Nancy\}"\} & d82a769d-\dots   \\
    \revision{Beijing}        & \{"\{Ellen,\revision{Jing}\}"\}                                         & 1a407bcb-\dots   \\
		\bottomrule
	\end{tabular}
\end{table}


We explain the query interface for provenance computation in GProM and
ProvSQL, and for probabilistic query evaluation in MayBMS and ProvSQL.
This helps understanding how these systems represent provenance and
uncertain information, and how to produce a benchmark adapted to each
system.

We consider again from Example~\ref{ex:semantics}
the \emph{Personnel} instance introduced in
Table \ref{tab:personnel} and the query $\qex$, which can be written in
SQL as follows:
\begin{minted}{sql}
  SELECT DISTINCT p1.city
  FROM personnel p1 JOIN personnel p2
  ON p1.city=p2.city AND p1.id<p2.id
\end{minted}

\paragraph*{Provenance computation}
To find the provenance of this query in GProM we need to use the
\mintkeyword{PROVENANCE OF}
construct in the query statement:
\begin{minted}[escapeinside=||]{sql}
  |\PYG{k}{PROVENANCE OF}| (SELECT DISTINCT p1.city
    FROM personnel p1 JOIN personnel p2
    ON p1.city=p2.city AND p1.id<p2.id)
\end{minted}
For this query, GProM produces the provenance derivation in Table~\ref{tab:gprom_query2}.
GProM attaches multiple additional attributes to the output tuples of a query.
In GProM's result, ``Paris'' appears three times: once due to
{\revision{David}} and
{\revision{Aaheli}}, once due to {\revision{Aaeheli}} and {Nancy}, and once due to
  {\revision{David}}
and {Nancy}. Each occurrence of {Paris} in GProM's output is accompanied
by the provenance attributes of the two input tuples that caused it to
appear. We see in particular that GProM disregards the
\mintkeyword{DISTINCT} keyword: the output of the query would be the
same without it.

To find similar provenance information for the same query using
ProvSQL, we first make the table provenance-aware by using the
\mintinline{sql}{add_provenance} UDF and then specialize the generic
provenance token, using a provenance mapping \emph{personnel\_name} where each tuple is
represented by the \textit{name} attribute, to why-provenance:
\begin{minted}{sql}
  SELECT p1.city, why(provenance(), 'personnel_name')
  FROM personnel p1 JOIN personnel p2
  ON p1.city=p2.city AND p1.id<p2.id
  GROUP BY p1.city
\end{minted}
Provenance evaluation is seen as a form of aggregation, which explains
why we transformed \mintkeyword{DISTINCT} into a \mintkeyword{GROUP BY}.
This results in a more compact representation, as shown in Table~\ref{tab:provsql_query2}.
ProvSQL computes the why-provenance for each output tuple as sets of
sets. Of course, in ProvSQL, any (m-)semiring can be used for the
computation, not just why-provenance: ProvSQL defines a number of common
examples, and a user can add new UDFs defining the operations of the
semiring of choice.

\begin{table}
	\centering
	\caption{The \emph{Personnel\_prob} table in MayBMS}\label{tab:personnel_p}
	\begin{tabular}{lllllll}
		\toprule
		\textbf{id} & \textbf{name} & \textbf{position} & \textbf{city} & \textbf{\_v0} & \textbf{\_d0} & \textbf{\_p0} \\
		\midrule
    1           & \revision{Juma}          & Director          & \revision{Nairobi}      & 1008          & 1             & 0.5           \\
		2           & Paul          & Janitor           & \revision{Nairobi}      & 1009          & 1             & 0.7           \\
    3           & \revision{David}          & Analyst           & Paris         & 10010         & 1             & 0.3           \\
		4           & Ellen         & Field agent       & \revision{Beijing}        & 10011         & 1             & 0.2           \\
    5           & \revision{Aaheli}     & Double agent      & Paris         & 10012         & 1             & 1.0           \\
		6           & Nancy         & HR                & Paris         & 10013         & 1             & 0.8           \\
    7           & \revision{Jing}         & Analyst           & \revision{Beijing}        & 10014         & 1             & 0.2           \\
		\bottomrule
	\end{tabular}
\end{table}
\begin{table}
	\centering
	\caption{Result of probability computation in ProvSQL (right) and
		MayBMS (left)}\label{tab:maybms_prob_eval}
	\begin{tabular}{ll}
		\toprule
		\textbf{city} & \textbf{conf} \\
		\midrule
		\revision{Nairobi}      & $0.35$        \\
		Paris         & $0.86$        \\
		\revision{Beijing}        & $0.04$        \\
		\bottomrule
	\end{tabular}
	\qquad
	\begin{tabular}{lll}
		\toprule
		\textbf{city} & \textbf{prob} & \textbf{provsql} \\
		\midrule
		\revision{Nairobi}      & 0.35          & d1b22232-\dots   \\
		Paris         & 0.86          & d82a769d-\dots   \\
		\revision{Beijing}        & 0.04          & 1a407bcb-\dots   \\
		\bottomrule
	\end{tabular}
\end{table}
\paragraph*{Probability computation}
We now consider the
TID relation introduced in Example~\ref{ex:probability}.

To compute probabilities in MayBMS, the following steps have to be
followed. We first use \mintkeyword{PICK TUPLES} to create a new
TID table with  probabilities given by some SQL expression. In this new table, named \emph{Personnel\_prob}, MayBMS
adds three extra attribute: an event variable \emph{\_v0}, a decision
value \emph{\_d0} and a probability \emph{\_p0}, as shown in
Table~\ref{tab:personnel_p}. In short, each unique event variable translate to a
set of mutually exclusive events, independent across event variables;
the decision value specifies the event, and the probability is that of
the event value. This is an encoding of the compact U-relation model for
probabilistic databases~\cite{DBLP:conf/icde/AntovaJKO08}, itself very
similar to that of probabilistic c-tables~\cite{greent06}.
MayBMS offers two functions for probabilistic qurery evaluation:
\mintinline{sql}{conf()} and \mintinline{sql}{tconf()} for computing the exact
probabilities (along with some functions for computing approximations).
The difference between these two functions is that
\mintinline{sql}{conf()} computes a probability of distinct results
(along with a \mintinline{sql}{GROUP BY} clause), while
\mintinline{sql}{tconf()} performs no duplicate elimination.
We can then use the following query:
\begin{minted}{sql}
  SELECT city, conf() FROM (
    SELECT p1.city
    FROM personnel_prob p1, personnel_prob p2
    WHERE p1.city=p2.city AND p1.id<p2.id) temp
  GROUP BY city
\end{minted}
whose result in MayBMS is given in Table~\ref{tab:maybms_prob_eval}
(left), matching what we computed in Example~\ref{ex:probability}.
Note that MayBMS does not support the \mintkeyword{JOIN} keyword
for queries with self-join, so the query needs to be rewritten with a
subquery as shown.

In ProvSQL, we first attach a probability on Boolean variables using the
\mintinline{sql}{set_prob()} UDF then, at query time, we use the
a UDF to perform probabilistic
query evaluation:
\begin{minted}{sql}
  SELECT p1.city, probability_evaluate(provenance())
  FROM personnel p1 JOIN personnel p2
  ON p1.city=p2.city AND p1.id<p2.id
  GROUP BY p1.city
\end{minted}
The result is shown in Table~\ref{tab:maybms_prob_eval} (right)
-- we
notice the result is still a provenance-aware table, with its
\emph{provsql} attribute.
\end{toappendix}
